\chapter{Release-quality discipline: tests, typecheck, lint, builds}
\label{ch:release-quality}

This chapter is about the habits that make a codebase shippable.
Relax is a teaching framework, but it also treats its test suite as an executable spec and its build pipeline as part of correctness.

The goal of this chapter is not to list commands---it's to explain \emph{why} each quality gate exists, what it catches, and how to debug failures without guessing.

If you only remember one thing: each gate is a different lens on correctness.
Typecheck defends your \emph{data model}, lint defends your \emph{human factors} (consistency and obvious mistakes), tests defend your \emph{observable behavior}, and builds defend your \emph{published artifact}.

\section{The quality-gates contract}

A ``release-quality'' change in this repo is expected to be green on:

\begin{itemize}
\item \textbf{Typecheck}: TypeScript compilation with \texttt{--noEmit}
\item \textbf{Lint}: ESLint rules for consistency and obvious mistakes
\item \textbf{Unit/integration tests}: Vitest (jsdom)
\item \textbf{Build}: Rollup bundles for ESM + CJS, plus type bundle generation on publish
\end{itemize}

\section*{Diagram: quality gates as a pipeline}
\begin{center}
\begin{tikzpicture}[
	node distance=8mm and 10mm,
	box/.style={draw, rounded corners, align=center, inner sep=7pt, fill=RelaxLight},
	arrow/.style={-Latex, thick}
]
\node[box] (change) {Change\\(code + docs)};
\node[box, right=of change] (type) {Typecheck\\\texttt{tsc --noEmit}};
\node[box, right=of type] (lint) {Lint\\\texttt{eslint}};
\node[box, right=of lint] (test) {Tests\\\texttt{vitest}};
\node[box, right=of test] (build) {Build\\Rollup + types};
\node[box, right=of build] (ship) {Publish\\(prepack)};

\draw[arrow] (change) -- (type);
\draw[arrow] (type) -- (lint);
\draw[arrow] (lint) -- (test);
\draw[arrow] (test) -- (build);
\draw[arrow] (build) -- (ship);
\end{tikzpicture}
\end{center}

This diagram is intentionally linear. In practice you often run these in parallel locally, but conceptually each step answers a different question: ``Do the types hold?'', ``Is the diff readable?'', ``Is the behavior correct?'', and ``Is the published artifact correct?''

These are not redundant; they fail on different classes of bugs.

\subsection{A tiny contract (inputs, outputs, failure modes)}

It helps to think about the pipeline as a function that takes a working tree and returns publishable artifacts:

\begin{itemize}
\item \textbf{Input}: TypeScript source under \texttt{src/}, plus experimental surfaces under \texttt{runtime/}, \texttt{compiler/}, and the TSX runtime under \texttt{relax-jsx/}.
\item \textbf{Output}: compiled JS bundles in \texttt{dist/esm/} and \texttt{dist/cjs/} and a declaration bundle in \texttt{dist-types/}.
\item \textbf{Success}: consumers can import from the package \texttt{exports} map and get the runtime semantics the tests describe.
\item \textbf{Failure modes}: ``works in tests'' but fails to typecheck; ``works locally'' but fails to bundle; ``bundles'' but exports the wrong files.
\end{itemize}

\section{Typecheck: make illegal states unrepresentable}

The repo's primary TypeScript config is \texttt{tsconfig.json}.
A few settings worth noticing:

\begin{itemize}
\item \texttt{strict: true} --- the baseline.
\item \texttt{noUncheckedIndexedAccess: true} --- forces you to handle missing keys.
\item \texttt{exactOptionalPropertyTypes: true} --- optional means ``may be absent'', not ``may be undefined''.
\item \texttt{jsx: react-jsx} and \texttt{jsxImportSource: relax-jsx} --- enables the TSX layer.
\end{itemize}

Two other settings are quietly important for how the codebase is built and consumed:

\begin{itemize}
\item \texttt{moduleResolution: Bundler}: TypeScript resolves imports similarly to modern bundlers.
In practice, this aligns the typechecker with how Rollup (and consumer bundlers) will interpret ESM.
\item \texttt{paths} for \texttt{relax-jsx}: TS can typecheck TSX code that imports the JSX runtime via the virtual specifier \texttt{relax-jsx}.
This is mirrored on the test side by the Vitest alias for \texttt{relax-jsx}.
\end{itemize}

The \texttt{typecheck} script (in \texttt{package.json}) runs:

\begin{quote}
\texttt{tsc -p tsconfig.json --noEmit}
\end{quote}

\section*{Diagram: why multiple tsconfig files exist}
\begin{center}
\begin{tikzpicture}[
	node distance=9mm and 12mm,
	box/.style={draw, rounded corners, align=left, inner sep=6pt, fill=RelaxLight},
	arrow/.style={-Latex, thick}
]
\node[box] (base) {\texttt{tsconfig.json}\\Developer truth\\strictness + DOM libs};
\node[box, below left=of base] (types) {\texttt{tsconfig.build-types.json}\\Emit \texttt{.d.ts} only\\to \texttt{dist-types/}};
\node[box, below right=of base] (rollup) {\texttt{tsconfig.rollup*.json}\\No declarations\\build JS for bundling};

\draw[arrow] (base) -- (types);
\draw[arrow] (base) -- (rollup);
\end{tikzpicture}
\end{center}

\subsection{Common typecheck failures (and how to debug)}

\begin{itemize}
\item \textbf{Indexing errors}: caused by \texttt{noUncheckedIndexedAccess}. Fix by narrowing or checking bounds.
\item \textbf{Optional property mismatches}: caused by \texttt{exactOptionalPropertyTypes}. Fix by removing explicit \texttt{undefined} assignments or changing types.
\item \textbf{DOM typing issues in tests}: Vitest uses jsdom; make sure \texttt{lib: ["DOM"]} is present (it is).
\end{itemize}

\subsection{Why there are multiple \texttt{tsconfig} files}

This repo uses a small family of configs to keep responsibilities separate:

\begin{itemize}
\item \texttt{tsconfig.json}: the ``truth'' for developer typechecking.
It enables declarations and maps, targets modern JavaScript (\texttt{ES2022}), and includes the DOM library.
\item \texttt{tsconfig.build-types.json}: extends the base config but turns on \emph{declaration-only} emit.
This produces the publishable \texttt{.d.ts} surface under \texttt{dist-types/}.
\item \texttt{tsconfig.rollup.json} and \texttt{tsconfig.rollup.cjs.json}: extend the base config but disable declaration emission.
These exist to keep Rollup builds focused on JavaScript output (ESM to \texttt{dist/esm/}, CJS to \texttt{dist/cjs/}) without creating duplicate \texttt{.d.ts} trees.
\end{itemize}

This split is release discipline in practice: build tooling stays fast and predictable, while types remain a first-class published artifact.

\section{Lint: keep the surface area consistent}

Lint config lives in \texttt{eslint.config.mjs}.
It uses the ESLint 9 flat config model and \texttt{@eslint/js} recommended rules.

The lint scripts are:

\begin{itemize}
\item \texttt{npm run lint} --- runs eslint on \texttt{src}
\item \texttt{npm run lint:fix} --- applies safe auto-fixes
\end{itemize}

Note: linting only \texttt{src} is a choice. HRBR runtime and compiler code are primarily guarded by tests and typecheck, but could be added to lint scope later.

\subsection{What lint catches that typecheck doesn't}

TypeScript proves things about types, but lint is often better at catching:

\begin{itemize}
\item accidental unused variables or imports
\item inconsistent patterns that make reviews harder
\item ``looks correct'' code that has a known foot-gun shape
\end{itemize}

Even with a mostly-recommended ESLint config, getting a clean lint run is part of keeping diffs small and reviewable.

\section{Tests: executable specs with jsdom}

All tests use Vitest (see \texttt{vitest.config.js}).
The environment is \texttt{jsdom}, which means:

\begin{itemize}
\item DOM APIs are available (document, elements, events)
\item you can test DOM manipulation deterministically without a real browser
\end{itemize}

Test inclusion patterns:

\begin{itemize}
\item \texttt{**/\_\_tests\_\_/**/*.test.ts}
\item \texttt{**/\_\_tests\_\_/**/*.test.tsx}
\end{itemize}

\subsection{The test suite structure}

This repo uses tests in three roles:

\begin{itemize}
\item \textbf{Unit tests}: small functions like diffs and helpers.
\item \textbf{Integration tests}: end-to-end flows (mount $\to$ patch $\to$ destroy, SSR $\to$ hydrate, compiler $\to$ mount).
\item \textbf{Property/fuzz tests}: seeded tests that exercise many update sequences (e.g. reconciler and signals).
\end{itemize}

If you change invariants (like reconciliation semantics), expect to update both unit tests and fuzz oracles.

\subsection{Tests-as-spec example: TSX wiring}

The TSX layer is a good example of how ``typecheck + tests'' combine into a stronger guarantee.
The typechecker is configured to understand TSX (\texttt{jsx: react-jsx} and \texttt{jsxImportSource: relax-jsx}), but tests still need to verify the \emph{runtime} behavior.

In \texttt{src/\_\_tests\_\_/jsx-tsx.test.ts}, the test sets up a host element in jsdom, runs an example, and asserts on the resulting HTML:

\begin{itemize}
\item it writes \texttt{<div id="app"></div>} into \texttt{document.body}
\item it locates the host element and passes it into \texttt{runHelloExample}
\item it asserts that the returned HTML matches \texttt{<div class="hello">Hello Relax</div>}
\end{itemize}

The important release-quality point is the scope of the guarantee:
this test doesn't just say ``TSX compiles''---it says ``the TSX runtime integrates with the DOM mounting layer and produces the expected structure''.

\subsection{Debugging failing tests (a workflow)}

A practical workflow:

\begin{enumerate}
\item Re-run only the failing file (Vitest supports file targeting; use the repo scripts).
\item Read the failure output and find which invariant changed.
\item Locate the corresponding runtime function and its other tests.
\item Add a minimal regression test that isolates the bug.
\item Only then, fix the implementation and re-run the focused test.
\end{enumerate}

When failures are timing-related (scheduler, signals), prefer deterministic schedulers and explicit ticking (the repo already uses this pattern).

\subsection{What to do when a test failure seems ``impossible''}

The fastest way to debug tough failures is to make the implicit invariant explicit:

\begin{enumerate}
\item Find the smallest assertion that fails.
\item Name the invariant in one sentence (e.g. ``patch preserves node identity when keys and tag match'').
\item Search for other tests that mention the same concept.
\item Read the corresponding implementation with the invariant in mind.
\end{enumerate}

This avoids the common trap of chasing symptoms across unrelated modules.

\section{Builds: ESM, CJS, and types}

Relax ships multiple entrypoints (see \texttt{package.json} \texttt{exports}):

\begin{itemize}
\item \texttt{"."} --- the VDOM library
\item \texttt{"./hrbr"} --- experimental HRBR runtime
\item \texttt{"./compiler"} --- experimental compiler
\end{itemize}

Rollup builds:

\begin{itemize}
\item ESM: \texttt{rollup.config.mjs} -> \texttt{dist/esm/}
\item CJS: \texttt{rollup.config.cjs.mjs} -> \texttt{dist/cjs/}
\end{itemize}

Type declaration builds:

\begin{itemize}
\item \texttt{npm run build:types} uses \texttt{tsconfig.build-types.json} (emit declarations only)
\item \texttt{npm run prepack} runs both JS builds and types before publishing
\end{itemize}

\subsection{Why builds are a quality gate}

A green test suite doesn't guarantee:

\begin{itemize}
\item exports are correct
\item bundling doesn't pull in unexpected dependencies
\item generated output is structurally valid for consumers
\end{itemize}

For example, the compiler bundle is large; build output sizes help you notice unexpected growth.

\subsection{Build discipline: match \texttt{exports} to outputs}

From \texttt{package.json}, Relax publishes multiple entry points (\texttt{"."}, \texttt{"./hrbr"}, \texttt{"./compiler"}).
Release discipline means each export must be correct in three dimensions:

\begin{itemize}
\item \textbf{ESM path} points at something in \texttt{dist/esm/}
\item \textbf{CJS path} points at something in \texttt{dist/cjs/} (with \texttt{.cjs} entry file names)
\item \textbf{Types path} points at something in \texttt{dist-types/}
\end{itemize}

It's easy to accidentally break this mapping when renaming files, changing entry points, or restructuring module boundaries.
That is why \texttt{prepack} builds both JS bundles and types before publish.

\section*{Diagram: package exports must match ESM, CJS, and types}
\begin{center}
\begin{tikzpicture}[
	node distance=9mm and 14mm,
	box/.style={draw, rounded corners, align=center, inner sep=6pt, fill=RelaxLight},
	arrow/.style={-Latex, thick}
]
\node[box] (exports) {\texttt{package.json exports}\\\texttt{"."} / \texttt{"./hrbr"} / \texttt{"./compiler"}};

\node[box, below left=of exports] (esm) {ESM\\\texttt{dist/esm/}};
\node[box, below=of exports] (cjs) {CJS\\\texttt{dist/cjs/}};
\node[box, below right=of exports] (types) {Types\\\texttt{dist-types/}};

\draw[arrow] (exports) -- (esm);
\draw[arrow] (exports) -- (cjs);
\draw[arrow] (exports) -- (types);
\end{tikzpicture}
\end{center}

\section{Release checklist (practical)}

A minimal release checklist aligned with this repo:

\begin{enumerate}
\item \textbf{Typecheck}: ensure TS still models your invariants.
\item \textbf{Lint}: ensure code style and basic safety checks pass.
\item \textbf{Tests}: run the suite; if you're changing HRBR internals, keep an eye on fuzz determinism.
\item \textbf{Build}: confirm ESM + CJS outputs compile.
\item \textbf{Bench smoke (optional but recommended)}: \texttt{build:bench:smoke} to catch catastrophic perf regressions.
\end{enumerate}

\section{Evidence from this repo}

This chapter is grounded in real scripts and configs:

\begin{itemize}
\item \texttt{package.json} scripts: \texttt{typecheck}, \texttt{lint}, \texttt{test}, \texttt{build}, \texttt{prepack}
\item \texttt{vitest.config.js} (jsdom environment)
\item \texttt{eslint.config.mjs} (flat config)
\item Rollup configs for bundling
\item TypeScript configs for declaration and rollup builds
\end{itemize}

In the next section of the book (appendices), we'll turn these practices into a learner-friendly ``build it from scratch'' checklist and a test-reading guide.

\section{Online resources}

These references mirror the key moving parts in this chapter:

\begin{itemize}
\item TypeScript \texttt{tsconfig} reference (compiler options, strictness flags):
\href{https://www.typescriptlang.org/tsconfig}{https://www.typescriptlang.org/tsconfig}
\item \texttt{exactOptionalPropertyTypes} (why ``optional'' is not the same as ``possibly undefined''):
\href{https://www.typescriptlang.org/tsconfig#exactOptionalPropertyTypes}{https://www.typescriptlang.org/tsconfig\#exactOptionalPropertyTypes}
\item \texttt{noUncheckedIndexedAccess} (making missing keys explicit):
\href{https://www.typescriptlang.org/tsconfig#noUncheckedIndexedAccess}{https://www.typescriptlang.org/tsconfig\#noUncheckedIndexedAccess}
\item TypeScript JSX / \texttt{jsxImportSource}:
\href{https://www.typescriptlang.org/docs/handbook/jsx.html}{https://www.typescriptlang.org/docs/handbook/jsx.html}
\item ESLint flat config (ESLint 9+):
\href{https://eslint.org/docs/latest/use/configure/configuration-files}{https://eslint.org/docs/latest/use/configure/configuration-files}
\item Vitest docs (jsdom environment, include patterns, reporters):
\href{https://vitest.dev/guide/}{https://vitest.dev/guide/}
\item jsdom project (what it simulates and what it doesn't):
\href{https://github.com/jsdom/jsdom}{https://github.com/jsdom/jsdom}
\item Rollup configuration and multi-entry builds:
\href{https://rollupjs.org/configuration-options/}{https://rollupjs.org/configuration-options/}
\item Node.js \texttt{package.json} \texttt{exports} field (how entry points are resolved):
\href{https://nodejs.org/api/packages.html#exports}{https://nodejs.org/api/packages.html\#exports}
\end{itemize}
